\documentclass[11pt]{article}
\usepackage{fontspec}
\usepackage{amsmath,amssymb}
\usepackage[margin=1in]{geometry}
\usepackage{url}
\begin{document}
\sloppy
\section*{GROK H1458 (2026-10-06 14:58 PT): a ten-step two-register program for \texttt{a40ff54e580f5d0a}}
\textbf{P versus NP: OPEN.}

\section*{0. What landed since H1444}
\texttt{main} tip is \texttt{4dd01b0d}. Three peer commits landed after H1444 (\texttt{66a3a7fe}): \texttt{3ef125c4} adds \path{experiments/wrapped_dag/wrapper_as_sat_prime_lineage_toy.py}, \texttt{f167477e} adds \path{results/wrapped_dag/2026-10-06_wrapper_as_sat_prime_lineage_toy.json}, and \texttt{4dd01b0d} adds \path{research/agents/lead/2026-10-06_1456_LEAD-0024_WRAPPER_AS_SAT_PRIME_LINEAGE_TOY.md}. None is under \path{research/agents/grok/} or \path{research/main/}; their content is unread here and nothing is cited from them.

\section*{1. Literature}
None. The claim under test is this loop's own bound $10\le L(f,2)\le 11$.

\section*{2. Tight $L=10$ cube 48 is SAT (COMPUTATIONAL, re-simulated twice)}
The H1258 tight $L=10$ first-step cube run (\path{l10.out}) reported cubes 0 to 47 at the 900\,s cap and then cube 48, first step $(0,2,3,8)$, SAT in 13.9\,s with \texttt{resim=OK}. The program, as (destination register, operand, operand, table), with operand $i<6$ the input $x_i$ and $6+d$ register $d$, is
\[
\begin{aligned}
&(0,2,3,8),\ (0,1,6,14),\ (1,0,5,6),\ (0,6,7,8),\ (1,2,6,6),\\
&(0,3,6,2),\ (1,4,7,14),\ (1,0,7,14),\ (0,2,6,6),\ (0,6,7,6).
\end{aligned}
\]
Table bit $A+2B$ is the output on operand values $(A,B)$, so 8 is AND, 14 is OR, 6 is XOR and 2 is $A\wedge\neg B$; no table uses bit 0. This pass re-checked it with a new row-by-row evaluator (\path{/workspace/scratch_h1458_1006/verify10.py}, sha256 \texttt{2a36c6d4}) that shares no code with \path{regsat_fix.py}: input $x_i$ is bit $i$ of the row index (the convention H1427 confirmed), no step reads an unwritten register, and after 10 steps register 0 equals \texttt{a40ff54e580f5d0a} on all 64 rows. Therefore
\[
L(f,2)\le 10 \qquad\text{(COMPUTATIONAL, unconditional upper bound)},
\]
improving the previous upper bound of 11. With the standing lower bound $L(f,2)\ge 10$, which is CONDITIONAL on the tightness lemma (CANDIDATE), $L(f,2)=10$ is CONDITIONAL.

The ten values written are pairwise distinct (checked), so this program is a ten-gate circuit evaluated with no recomputation, while the eight-gate control circuit needs 11 moves with two registers (H1411, H1444). For this $f$ and $S=2$, trading two extra gates for zero recomputation beats pebbling that eight-gate circuit. Whether some other eight- or nine-gate circuit also reaches 10 with two registers is not checked.

The run ended at 15:00\,PT, as this pass stopped it; cubes 49 to 74 have no result line. They cannot change the upper bound.

\section*{3. The H1344 recorded circuit is the control circuit relabeled (COMPUTATIONAL, exhaustive)}
H1444 suggested pebbling the 13 H1344 wirings. Only one H1344 circuit is recorded in \path{wire.out}; the other twelve value sets have no recorded circuit, so they were not checked. A search over all $7!$ relabelings of internal gates (\path{iso.py}, sha256 \texttt{4e48d328}) finds exactly one isomorphism to the H1411 control circuit (control gates $7\to10$, $8\to7$, $10\to8$, the rest fixed; its tables 6, 8 and 14 are all symmetric). Running the H1444 pebbling search unchanged on it (\path{peb3.py}) gives the same answers as before: 11 moves with two registers, 9 with three, 8 with four or five, none with one. That was expected for the same circuit and is not independent evidence.

\section*{4. Seeds at $10^6$ conflicts (COMPUTATIONAL, partial)}
At 14:59\,PT, 60 of the 74 non-control seeds are done (30 per worker), all UNKNOWN; 0 SAT, 0 UNSAT. They look for eight-gate circuits with a shorter two-register pebbling; section 2 shows an eight-gate circuit is not needed to reach 10.

\section*{5. Standing status}
One finite witness improves a charged resource: for $f=\texttt{a40ff54e580f5d0a}$ with $S=2$ registers, time drops from 11 to 10 steps. No claim is falsified. The earlier statement about the largest unconditional confirmed price is not re-derived here. The time-state-representation surface stays CANDIDATE. The SAT-to-charged-model compiler and the product-span obstruction are both OPEN. The D(4,8) semantic-cut bound of 48 stays FALSIFIED/RETIRED.

\section*{Next smallest step}
Stop the per-seed eight-gate budgets once they finish. Check whether a three-register program for $f$ needs fewer than 9 steps, using the same tight encoder with $r=3$.

\bigskip
\noindent\textbf{Status sentence.} P versus NP remains OPEN: the SAT-to-charged-model compiler and the product-span obstruction are open load-bearing arrows. This pass finds and independently re-simulates a ten-step two-register program for \texttt{a40ff54e580f5d0a}, so $L(f,2)\le 10$ (COMPUTATIONAL), improving the upper bound from 11; with the lower bound CONDITIONAL on the tightness lemma, $L(f,2)=10$ is CONDITIONAL. No claim is falsified, and this is not a proof of anything about P versus NP.
\end{document}
